%% AMS-LaTeX Created with the Wolfram Language : www.wolfram.com

\documentclass[12pt]{article}

\usepackage[letterpaper,margin=1in]{geometry}
\usepackage{amsmath,amssymb,setspace}
\usepackage{graphicx}
\usepackage{bm}
\usepackage{booktabs}
\usepackage{multicol,multirow}
\graphicspath{{./figs/}}
\usepackage{tcolorbox}
\usepackage[framed,numbered,autolinebreaks,useliterate]{mcode}
\usepackage{enumitem}
\usepackage{xcolor}
\usepackage{listings}
\usepackage{hyperref}
\usepackage[cal=boondox]{mathalfa}

% Define colors for lstlisting
\definecolor{codegreen}{rgb}{0,0.6,0}
\definecolor{codegray}{rgb}{0.5,0.5,0.5}
\definecolor{codepurple}{rgb}{0.58,0,0.82}
\definecolor{backcolour}{rgb}{0.95,0.95,0.95}

% First, add this in your preamble if not already done
\usepackage{setspace}

% Then, modify your lstlisting style
\lstdefinestyle{mystyle}{
    backgroundcolor=\color{backcolour},   
    commentstyle=\color{codegreen},
    keywordstyle=\color{magenta},
    numberstyle=\tiny\color{codegray},
    stringstyle=\color{codepurple},
    basicstyle=\ttfamily\normalsize\setstretch{1}, % set single spacing explicitly here
    breakatwhitespace=false,         
    breaklines=true,                 
    captionpos=b,                    
    keepspaces=true,                 
    numbers=none,  
    columns=flexible,         
    numbersep=5pt,                  
    showspaces=false,                
    showstringspaces=false,
    showtabs=false,                  
    tabsize=2,
    literate={'}{{\textquotesingle}}1 {''}{{\textquotedbl}}2
}

\lstset{style=mystyle}

\hypersetup{
    colorlinks=true,
    linkcolor=blue,
    filecolor=magenta,      
    urlcolor=cyan,
    pdftitle={Your Document Title},
    pdfpagemode=FullScreen,
}

\def\mathbi#1{\textbf{\em #1}}
\newcommand{\fullunderline}{\noindent\makebox[\linewidth]{\rule{\linewidth}{0.1pt}}\\[-2em]}

\begin{document}
\doublespacing	

\title{CE 401/5501 Homework 08\\[2in]
\textbf{Name: \rule{2in}{0.15mm}}}
\author{}
\date{}
\maketitle

\newcommand*{\I}{\imath} % imaginary unit i

\newpage

\section*{Concept Questions}

\noindent \textbf{Question 1}

First, review the SVM model (Topic 06 notes) and the perceptron model (Topic 05) concepts. If the perceptron uses a hard activation function (a sign function), then a linear SVM and a perceptron both have the following functional form:
\[
m(\bm{x}) = f\bigl[u(\bm{x})\bigr] = \mathrm{sign}\bigl[u(\bm{x})\bigr]
\]
with
\[
u(\bm{x}) = \langle \bm{w}, \bm{x} \rangle + w_0,
\]
where $\bm{w}$ is the weight vector and $w_0$ is the bias (intercept) parameter.

Discuss by listing a few points that differentiate a linear SVM from a perceptron in the binary classification case.

\newpage
\noindent \textbf{Question 2}

Given a linear SVM, the loss function is defined as 
\[
h_i = \max\bigl[0,\, 1- u(\bm{x})\times y_i\bigr],
\]
where $u(\bm{x})$ is the decision function and $y_i$ is the true label of the data point $\bm{x}_i$. 

\begin{itemize}
    \item Plot manually the loss function $h_i(u)$ versus $u$. \textit{Hint: Discuss the two cases $y_i = 1$ and $y_i = -1$, with $u$ varying from, for example, $-3$ to $3$. You can first write it as a piece-wise function of $u$.}
    
    \item Then, summarize how the loss changes when $|u| \geq 1$ and when $|u| \leq 1$.
\end{itemize}

\newpage
\noindent \textbf{Question 3}

When using a polynomial kernel SVM, its kernel function is expressed as
\begin{align*}
  K(\bm{x}, \bm{x}_i) & = \langle \Phi(\bm{x}), \Phi(\bm{x}_i)\rangle\\[1mm]
  & = \bigl(\langle \bm{x}, \bm{x}_i\rangle + 1\bigr)^d,
\end{align*}
where $d$ is the degree of the polynomial.

To achieve this kernel, the transformed feature $\Phi(\bm{x})$ is mapped to a higher-dimensional space. If $\bm{x} \in \mathbb{R}^D$, then $\Phi(\bm{x}) \in \mathbb{R}^{D_k}$ with 
\[
D_k = \binom{D+d}{d}.
\]

Given this:
\begin{itemize}
    \item If $D = 2$ and $d = 2$, namely, a 2nd-order polynomial kernel is used for features in a 2-dimensional space, calculate $D_k$ to determine the dimension of the transformed space.
    
    \item Use ChatGPT or Gemini and prompt something like: 
    \begin{quote}
    ``In a polynomial kernel SVM, if the input feature is $\bm{x} = [x_1, x_2]^T$, and the kernel is $k(\bm{x}, \bm{x}') = (\langle \bm{x}, \bm{x}'\rangle + 1)^2$, i.e., using a 2nd-order polynomial kernel, what does the transformation mapping look like?''
    \end{quote}
    Report what you find and write out the transformed feature of $\bm{x} = [x_1, x_2]^T$:
    \[
    \Phi(\bm{x}) = ?
    \]
    
\end{itemize}

\newpage
\section*{Programming Practice on SVMs}

In this practice, we continue the class demonstration using the same \emph{concrete data} (please download the \href{https://umsystem.instructure.com/courses/280496/files/33240224?module_item_id=9853160}{class notebook file}) and further study the behavior of different SVM classifiers.

First, please change the train-test split ratio; in this exercise, we will use $50\%$ of the data as the testing set (as follows).
\begin{verbatim}
X_train, X_test, y_train, y_test = ...
train_test_split(X, y, test_size=0.5, random_state=107)
\end{verbatim}

The following tasks are to be completed:
\begin{itemize}
    \item Using a linear SVM (i.e., with \texttt{kernel='linear'} and $C=1$), report the confusion matrix and the overall accuracy on both the training and testing sets. Additionally, display the decision boundary and margin plots using the provided sample code.
    
    \item Using polynomial kernels, consider the following experimental matrix (where $d$ is the degree of the polynomial and $C$ is the regularization parameter):
    % Please add the following required packages to your document preamble:
    % \usepackage{booktabs}
    % \usepackage{multirow}
    \begin{table}[htb]
    \centering
    \begin{tabular}{|lc|ccc|}
    \hline
    \multicolumn{2}{|c|}{\multirow{2}{*}{\textit{OA}}} & \multicolumn{3}{c|}{$C$}                                 \\ \cline{3-5} 
    \multicolumn{2}{|c|}{}                             & \multicolumn{1}{c|}{0.1} & \multicolumn{1}{c|}{1} & 10 \\ \hline
    \multicolumn{1}{|l|}{\multirow{3}{*}{$d$}} & 2 & \multicolumn{1}{c|}{} & \multicolumn{1}{c|}{} &  \\ \cline{2-5} 
    \multicolumn{1}{|l|}{}             & 3             & \multicolumn{1}{c|}{}    & \multicolumn{1}{c|}{}  &    \\ \cline{2-5} 
    \multicolumn{1}{|l|}{}             & 5             & \multicolumn{1}{c|}{}    & \multicolumn{1}{c|}{}  &    \\ \hline
    \end{tabular}
    \end{table}
    
    That is, you will consider nine combinations of hyperparameters; for example, $(C = 0.1,\, d = 2)$, $(C = 1,\, d = 2)$, etc.
    
    \begin{itemize}
        \item For each combination of hyperparameters, report the overall accuracy (OA) on the testing data and fill in the table above.
        \item Compared with the benchmark result from the first task (using a linear kernel with $C=1$), indicate which hyperparameter setting yields the best accuracy.
        \item Based on the results, for a fixed $C$ (e.g., $C=1$ or $C=10$), does the accuracy of the model increase as $d$ increases? Also, for a fixed $d$, is there a consistent trend in model accuracy as $C$ varies from 0.1 to 10 across different $d$ values?
    \end{itemize}

\end{itemize}

\end{document}
